\documentclass[12pt,a4paper]{article}
\usepackage{../../../myconfig}

\begin{document}

% =========================================================================
% BÀI TẬP VỀ NHÀ (BTVN): BỘ ĐỀ TOÀN DIỆN 35 CÂU (KHÔNG TRÙNG LẶP BTDD)
% =========================================================================
\titlebox{Chuyên đề: Oxyz}{Bổ sung kiến thức Oxyz}{Oxyz: Bổ sung kiến thức}

\vspace{-0.25cm}

% Câu 1 (Câu 390 trong tài liệu - Góc giữa hai mặt phẳng)
\questionLinesManual{0}{1}{%
Trong không gian $Oxyz$, góc giữa hai mặt phẳng $(P)\colon 8x - 4y - 8z - 11 = 0$ và $(Q)\colon \sqrt{2}x - \sqrt{2}y + 7 = 0$ bằng
\begin{tasks}(4)
    \task $\dfrac{\pi}{4}$.
    \task $\dfrac{\pi}{2}$.
    \task $\dfrac{\pi}{6}$.
    \task $\dfrac{\pi}{3}$.
\end{tasks}
}{}

\vspace{0.15cm}

% Câu 2 (Câu 412 trong tài liệu - Góc/VTTĐ 2 đường thẳng vuông góc)
\questionLinesManual{0}{2}{%
Trong không gian $Oxyz$, cho hai đường thẳng $\Delta_1\colon \begin{cases} x = -3 + 2t \\ y = 1 - t \\ z = -1 + 4t \end{cases}$ và $\Delta_2\colon \dfrac{x + 4}{3} = \dfrac{y + 2}{2} = \dfrac{z - 4}{-1}$. Khẳng định nào sau đây là đúng?
\begin{tasks}(2)
    \task $\Delta_1$ và $\Delta_2$ chéo nhau và vuông góc nhau.
    \task $\Delta_1$ cắt và không vuông góc với $\Delta_2$.
    \task $\Delta_1$ cắt và vuông góc với $\Delta_2$.
    \task $\Delta_1$ và $\Delta_2$ song song với nhau.
\end{tasks}
}{}

\vspace{0.15cm}

% Câu 3 (Câu 432 trong tài liệu - Đường thẳng vuông góc mặt phẳng)
\questionLinesManual{0}{3}{%
Trong không gian với hệ tọa độ $Oxyz$, cho đường thẳng $d\colon \dfrac{2x - 2}{3n} = \dfrac{y + 1}{4} = \dfrac{3z + 6}{2m}$ ($m, n \neq 0$) và mặt phẳng $(P)\colon 3x + 4y - 2z + 5 = 0$. Khi đường thẳng $d$ vuông góc với mặt phẳng $(P)$ thì tổng $m + n$ bằng
\begin{tasks}(4)
    \task $1$.
    \task $-1$.
    \task $3$.
    \task $-5$.
\end{tasks}
}{}

\vspace{0.15cm}

% Câu 4 (Câu 472 trong tài liệu - Góc/VTTĐ 2 đường thẳng)
\questionLinesManual{0}{4}{%
Trong không gian với hệ tọa độ $Oxyz$, cho hai đường thẳng $d_1\colon \dfrac{x - 1}{2} = \dfrac{y - 7}{1} = \dfrac{z}{4}$ và $d_2\colon \dfrac{x + 1}{1} = \dfrac{y - 2}{2} = \dfrac{z - 2}{-1}$. Mệnh đề nào sau đây là đúng?
\begin{tasks}(2)
    \task $d_1$ và $d_2$ vuông góc với nhau và cắt nhau.
    \task $d_1$ và $d_2$ song song với nhau.
    \task $d_1$ và $d_2$ trùng nhau.
    \task $d_1$ và $d_2$ chéo nhau.
\end{tasks}
}{}

\vspace{0.15cm}

% Câu 5 (Câu 392 trong tài liệu - Khoảng cách từ điểm đến mặt phẳng)
\questionLinesManual{0}{5}{%
Trong không gian với hệ tọa độ $Oxyz$, cho mặt phẳng $(P)\colon x - 2y - 2z + 5 = 0$ và điểm $A(-1; 3; -2)$. Khoảng cách từ điểm $A$ đến mặt phẳng $(P)$ bằng
\begin{tasks}(4)
    \task $d = 1$.
    \task $d = \dfrac{2}{3}$.
    \task $d = \dfrac{3\sqrt{14}}{14}$.
    \task $d = \dfrac{\sqrt{14}}{7}$.
\end{tasks}
}{}

\vspace{0.15cm}

% Câu 6 (Câu 393 trong tài liệu - Khoảng cách từ điểm đến mặt phẳng)
\questionLinesManual{0}{6}{%
Trong không gian $Oxyz$, cho mặt phẳng $(P)\colon 2x - 3y + 6z + 19 = 0$ và điểm $A(-2; 4; 3)$. Gọi $d$ là khoảng cách từ $A$ đến mặt phẳng $(P)$. Khi đó $d$ bằng
\begin{tasks}(4)
    \task $d = 4$.
    \task $d = 2$.
    \task $d = 1$.
    \task $d = 3$.
\end{tasks}
}{}

\vspace{0.15cm}

% Câu 7 (Câu 394 trong tài liệu - Khoảng cách từ điểm đến mặt phẳng)
\questionLinesManual{0}{7}{%
Trong không gian với hệ tọa độ $Oxyz$, cho mặt phẳng $(P)\colon 2x + 2y - z + 3 = 0$ và điểm $M(1; -2; -1)$. Khi đó khoảng cách từ điểm $M$ đến mặt phẳng $(P)$ bằng
\begin{tasks}(4)
    \task $\dfrac{8}{3}$.
    \task $\dfrac{10}{3}$.
    \task $0$.
    \task $\dfrac{2}{3}$.
\end{tasks}
}{}

\vspace{0.15cm}

% Câu 8 (Câu 395 trong tài liệu - Khoảng cách từ điểm đến mặt phẳng)
\questionLinesManual{0}{8}{%
Trong không gian $Oxyz$, cho mặt phẳng $(P)\colon 2x - 2y - z + 3 = 0$ và điểm $M(1; -2; 13)$. Khoảng cách $d$ từ điểm $M$ đến $(P)$ bằng
\begin{tasks}(4)
    \task $d = \dfrac{4}{3}$.
    \task $d = \dfrac{7}{3}$.
    \task $d = \dfrac{10}{3}$.
    \task $d = -\dfrac{4}{3}$.
\end{tasks}
}{}

\vspace{0.15cm}

% Câu 9 (Câu 396 trong tài liệu - Khoảng cách từ điểm đến mặt phẳng)
\questionLinesManual{0}{9}{%
Trong không gian với hệ trục tọa độ $Oxyz$, cho mặt phẳng $(P)\colon x + 2y - 2z + 1 = 0$ và điểm $M(1; -2; 2)$. Tính khoảng cách từ điểm $M$ đến mặt phẳng $(P)$.
\begin{tasks}(4)
    \task $d(M, (P)) = 2$.
    \task $d(M, (P)) = \dfrac{2}{3}$.
    \task $d(M, (P)) = \dfrac{10}{3}$.
    \task $d(M, (P)) = 3$.
\end{tasks}
}{}

\vspace{0.15cm}

% Câu 10 (Câu 398 trong tài liệu - Khoảng cách từ điểm đến mặt phẳng)
\questionLinesManual{0}{10}{%
Trong không gian $Oxyz$, cho mặt phẳng $(P)\colon 6x - 3y + 2z - 6 = 0$. Tính khoảng cách $d$ từ điểm $M(1; -2; 3)$ đến mặt phẳng $(P)$.
\begin{tasks}(4)
    \task $d = \dfrac{12\sqrt{85}}{85}$.
    \task $d = \dfrac{12}{7}$.
    \task $d = \dfrac{\sqrt{31}}{7}$.
    \task $d = \dfrac{18}{7}$.
\end{tasks}
}{}

\vspace{0.15cm}

% Câu 11 (Câu 447 trong tài liệu - Khoảng cách từ điểm đến mặt phẳng)
\questionLinesManual{0}{11}{%
Trong không gian với hệ trục tọa độ $Oxyz$, cho điểm $M(1; 2; -3)$ và mặt phẳng $(P)\colon x - 2y + 2z + 3 = 0$. Khoảng cách từ $M$ đến mặt phẳng $(P)$ là
\begin{tasks}(4)
    \task $1$.
    \task $2$.
    \task $3$.
    \task $4$.
\end{tasks}
}{}

\vspace{0.15cm}

% Câu 12 (Câu 451 trong tài liệu - Khoảng cách từ điểm đến mặt phẳng)
\questionLinesManual{0}{12}{%
Trong không gian với hệ tọa độ $Oxyz$, tính khoảng cách từ điểm $M(1; 2; -3)$ đến mặt phẳng $(P)\colon x + 2y - 2z - 2 = 0$.
\begin{tasks}(4)
    \task $1$.
    \task $\dfrac{11}{3}$.
    \task $\dfrac{1}{3}$.
    \task $3$.
\end{tasks}
}{}

\vspace{0.15cm}

% Câu 13 (Câu 441 trong tài liệu - Khoảng cách từ điểm đến mặt phẳng)
\questionLinesManual{0}{13}{%
Trong không gian với hệ tọa độ $Oxyz$, cho bốn điểm $A(-1; 2; 1)$, $B(-4; 2; -2)$, $C(-1; -1; -2)$ và $D(-5; -5; 2)$. Tính khoảng cách từ điểm $D$ đến mặt phẳng $(ABC)$.
\begin{tasks}(4)
    \task $d = \sqrt{3}$.
    \task $d = 2\sqrt{3}$.
    \task $d = 3\sqrt{3}$.
    \task $d = 4\sqrt{3}$.
\end{tasks}
}{}

\vspace{0.15cm}

% Câu 14 (Câu 450 trong tài liệu - Khoảng cách từ điểm đến mặt phẳng)
\questionLinesManual{0}{14}{%
Trong không gian với hệ tọa độ $Oxyz$, cho ba điểm $A(1; 0; 2)$, $B(1; 1; 1)$ và $C(2; 3; 0)$. Tính khoảng cách $h$ từ gốc tọa độ $O$ đến mặt phẳng $(ABC)$.
\begin{tasks}(4)
    \task $h = \sqrt{3}$.
    \task $h = \dfrac{1}{3}$.
    \task $h = 3$.
    \task $h = \dfrac{\sqrt{3}}{3}$.
\end{tasks}
}{}

\vspace{0.15cm}

% Câu 15 (Câu 454 trong tài liệu - Khoảng cách từ điểm đến đường thẳng)
\questionLinesManual{0}{15}{%
Trong không gian với hệ tọa độ $Oxyz$, cho đường thẳng $d\colon \dfrac{x - 1}{1} = \dfrac{y - 2}{2} = \dfrac{z + 2}{-2}$. Tính khoảng cách từ điểm $M(-2; 1; -1)$ tới $d$.
\begin{tasks}(4)
    \task $\dfrac{5\sqrt{2}}{3}$.
    \task $\dfrac{5\sqrt{2}}{2}$.
    \task $\dfrac{\sqrt{2}}{3}$.
    \task $\dfrac{5}{3}$.
\end{tasks}
}{}

\vspace{0.15cm}

% Câu 16 (Câu 446 trong tài liệu - Khoảng cách giữa 2 mặt phẳng song song)
\questionLinesManual{0}{16}{%
Trong không gian với hệ trục $Oxyz$, khoảng cách giữa hai mặt phẳng song song $(\alpha)\colon x - 2y - 2z + 4 = 0$ và $(\beta)\colon -x + 2y + 2z - 7 = 0$ là
\begin{tasks}(4)
    \task $1$.
    \task $-1$.
    \task $3$.
    \task $0$.
\end{tasks}
}{}

\vspace{0.15cm}

% Câu 17 (Câu 381 trong tài liệu - Mặt phẳng song song)
\questionLinesManual{0}{17}{%
Trong không gian với hệ tọa độ $Oxyz$, cho mặt phẳng $(P)$ có phương trình $3x + 2y - 3z + 1 = 0$. Phát biểu nào sau đây là sai?
\begin{tasks}(1)
    \task Phương trình của mặt phẳng song song với $(P)$ là $3x + 2y - 3z + 2 = 0$.
    \task Phương trình của mặt phẳng song song với $(P)$ là $6x + 4y - 6z - 1 = 0$.
    \task Phương trình của mặt phẳng song song với $(P)$ là $-3x - 2y + 3z - 5 = 0$.
    \task Phương trình của mặt phẳng song song với $(P)$ là $-3x - 2y + 3z - 1 = 0$.
\end{tasks}
}{}

\vspace{0.15cm}

% Câu 18 (Câu 383 trong tài liệu - Vị trí tương đối 2 mặt phẳng)
\questionLinesManual{0}{18}{%
Trong không gian $Oxyz$, cho hai mặt phẳng $(P)\colon 2x - 3y + z - 4 = 0$ và $(Q)\colon 5x - 3y - 2z - 7 = 0$. Vị trí tương đối của hai mặt phẳng $(P)$ và $(Q)$ là
\begin{tasks}(2)
    \task Song song.
    \task Cắt nhưng không vuông góc.
    \task Vuông góc.
    \task Trùng nhau.
\end{tasks}
}{}

\vspace{0.15cm}

% Câu 19 (Câu 404 trong tài liệu - Vị trí tương đối các mặt phẳng)
\questionLinesManual{0}{19}{%
Trong không gian với hệ tọa độ $Oxyz$, cho ba mặt phẳng $(P)\colon 3x + y + z - 4 = 0$, $(Q)\colon 3x + y + z + 5 = 0$ và $(R)\colon 2x - 3y - 3z + 1 = 0$. Xét các mệnh đề: (1) $(P) // (Q)$; (2) $(P) \perp (R)$. Khẳng định nào sau đây đúng?
\begin{tasks}(2)
    \task (1) đúng, (2) sai.
    \task (1) sai, (2) đúng.
    \task (1) đúng, (2) đúng.
    \task (1) sai, (2) sai.
\end{tasks}
}{}

\vspace{0.15cm}

% Câu 20 (Câu 406 trong tài liệu - Vị trí tương đối 2 mặt phẳng)
\questionLinesManual{0}{20}{%
Trong không gian với hệ tọa độ $Oxyz$, cho hai mặt phẳng $(P)\colon -2x + 6y - 4z - 1 = 0$ và $(Q)\colon x - 3y - 2z + 1 = 0$. Mệnh đề nào sau đây là đúng?
\begin{tasks}(2)
    \task $(P)$ cắt và không vuông góc với $(Q)$.
    \task $(P)$ vuông góc với $(Q)$.
    \task $(P)$ song song với $(Q)$.
    \task $(P)$ và $(Q)$ trùng nhau.
\end{tasks}
}{}

\vspace{0.15cm}

% Câu 21 (Câu 407 trong tài liệu - Tham số để 2 mặt phẳng song song)
\questionLinesManual{0}{21}{%
Trong không gian với hệ tọa độ $Oxyz$, mặt phẳng $(P)\colon x + my + 3z + 2 = 0$ và mặt phẳng $(Q)\colon nx + y + z + 7 = 0$ song song với nhau khi
\begin{tasks}(2)
    \task $m = n = 1$.
    \task $m = 3;\ n = \dfrac{1}{3}$.
    \task $m = 2;\ n = \dfrac{1}{3}$.
    \task $m = 3;\ n = \dfrac{1}{2}$.
\end{tasks}
}{}

\vspace{0.15cm}

% Câu 22 (Câu 386 trong tài liệu - Tham số để d vuông góc P)
\questionLinesManual{0}{22}{%
Cho đường thẳng $d\colon \begin{cases} x = 2 - 3t \\ y = 5 + 7t \\ z = 4 + (m - 3)t \end{cases}$ và mặt phẳng $(P)\colon 3x - 7y + 13z - 91 = 0$. Tìm giá trị của tham số $m$ để $d$ vuông góc với $(P)$.
\begin{tasks}(4)
    \task $13$.
    \task $-10$.
    \task $-13$.
    \task $10$.
\end{tasks}
}{}

\vspace{0.15cm}

% Câu 23 (Câu 387 trong tài liệu - Tham số để d song song P)
\questionLinesManual{0}{23}{%
Trong không gian với hệ trục tọa độ $Oxyz$, đường thẳng $d\colon \dfrac{x - 1}{2} = \dfrac{y + 2}{-1} = \dfrac{z + 1}{1}$ song song với mặt phẳng $(P)\colon x + y - z + m = 0$ khi và chỉ khi
\begin{tasks}(4)
    \task $\forall m \in \mathbb{R}$.
    \task $m = 0$.
    \task $m \neq 0$.
    \task $m \neq 2$.
\end{tasks}
}{}

\vspace{0.15cm}

% Câu 24 (Câu 388 trong tài liệu - Vị trí tương đối d và P)
\questionLinesManual{0}{24}{%
Trong không gian với hệ tọa độ $Oxyz$, cho mặt phẳng $(P)\colon 2x - 5y - 3z - 7 = 0$ và đường thẳng $d\colon \dfrac{x - 2}{2} = \dfrac{y}{-1} = \dfrac{z + 1}{3}$. Kết luận nào dưới đây là đúng?
\begin{tasks}(4)
    \task $d // (P)$.
    \task $d$ cắt $(P)$.
    \task $d \perp (P)$.
    \task $(P)$ chứa $d$.
\end{tasks}
}{}

\vspace{0.15cm}

% Câu 25 (Câu 389 trong tài liệu - Vị trí tương đối d và P)
\questionLinesManual{0}{25}{%
Cho đường thẳng $d\colon \dfrac{x - 1}{1} = \dfrac{y - 1}{2} = \dfrac{z - 2}{-3}$ và mặt phẳng $(\alpha)\colon x + y + z - 4 = 0$. Trong các khẳng định sau, khẳng định nào đúng?
\begin{tasks}(4)
    \task $d \subset (\alpha)$.
    \task $d // (\alpha)$.
    \task $d \perp (\alpha)$.
    \task $d$ cắt $(\alpha)$.
\end{tasks}
}{}

\vspace{0.15cm}

% Câu 26 (Câu 425 trong tài liệu - Tham số để d song song P)
\questionLinesManual{0}{26}{%
Trong không gian với hệ tọa độ $Oxyz$, cho đường thẳng $d\colon \dfrac{x - 2}{1} = \dfrac{y - 1}{1} = \dfrac{z - 1}{-1}$ và mặt phẳng $(P)\colon x + my + (m^2 - 1)z - 7 = 0$ ($m$ là tham số thực). Tìm $m$ để $d$ song song với $(P)$.
\begin{tasks}(4)
    \task $\left[\begin{aligned} m = -1 \\ m = 2 \end{aligned}\right.$.
    \task $m = -1$.
    \task $m = 2$.
    \task $m = 1$.
\end{tasks}
}{}

\vspace{0.15cm}

% Câu 27 (Câu 427 trong tài liệu - Đường thẳng vuông góc mặt phẳng)
\questionLinesManual{0}{27}{%
Trong không gian với hệ tọa độ $Oxyz$, đường thẳng $\Delta\colon \dfrac{x}{1} = \dfrac{y}{1} = \dfrac{z}{2}$ vuông góc với mặt phẳng nào trong các mặt phẳng sau?
\begin{tasks}(2)
    \task $(P)\colon x + y + z = 0$.
    \task $(\beta)\colon x + y - z = 0$.
    \task $(\alpha)\colon x + y + 2z = 0$.
    \task $(Q)\colon x + y - 2z = 0$.
\end{tasks}
}{}

\vspace{0.15cm}

% Câu 28 (Câu 428 trong tài liệu - Vị trí tương đối d và P)
\questionLinesManual{0}{28}{%
Cho đường thẳng $d\colon \dfrac{x - 1}{-1} = \dfrac{y}{2} = \dfrac{z - 3}{4}$ và mặt phẳng $(P)\colon 2x - y + z - 5 = 0$. Xét vị trí tương đối của $d$ và $(P)$.
\begin{tasks}(2)
    \task $d$ nằm trên $(P)$.
    \task $d$ song song với $(P)$.
    \task $d$ cắt và không vuông góc với $(P)$.
    \task $d$ vuông góc với $(P)$.
\end{tasks}
}{}

\vspace{0.15cm}

% Câu 29 (Câu 431 trong tài liệu - Tìm giao điểm d và P)
\questionLinesManual{0}{29}{%
Trong không gian với hệ trục tọa độ $Oxyz$, cho đường thẳng $d\colon \dfrac{x - 2}{-3} = \dfrac{y}{1} = \dfrac{z + 1}{2}$ và mặt phẳng $(P)\colon x + 2y - 3z + 2 = 0$. Tọa độ giao điểm $M$ của $d$ và $(P)$ là
\begin{tasks}(4)
    \task $M(-1; 1; 1)$.
    \task $M(2; 0; -1)$.
    \task $M(1; 0; 1)$.
    \task $M(5; -1; -3)$.
\end{tasks}
}{}

\vspace{0.15cm}

% Câu 30 (Câu 434 trong tài liệu - Vị trí tương đối d và P)
\questionLinesManual{0}{30}{%
Trong không gian với hệ tọa độ $Oxyz$, cho đường thẳng $d\colon \dfrac{x + 1}{1} = \dfrac{y}{-3} = \dfrac{z - 5}{-1}$ và mặt phẳng $(P)\colon x - 3y + 2z + 6 = 0$. Mệnh đề nào dưới đây đúng?
\begin{tasks}(2)
    \task $d$ cắt và không vuông góc với $(P)$.
    \task $d$ vuông góc với $(P)$.
    \task $d$ song song với $(P)$.
    \task $d$ nằm trong $(P)$.
\end{tasks}
}{}

\vspace{0.15cm}

% Câu 31 (Câu 408 trong tài liệu - Vị trí tương đối 2 đường thẳng)
\questionLinesManual{0}{31}{%
Trong không gian với hệ tọa độ $Oxyz$, cho hai đường thẳng $d\colon \dfrac{x - 2}{2} = \dfrac{y + 4}{3} = \dfrac{1 - z}{-2}$ và $d'\colon \begin{cases} x = 4t \\ y = 1 + 6t \\ z = -1 + 4t \end{cases}$. Vị trí tương đối giữa hai đường thẳng $d$ và $d'$ là
\begin{tasks}(2)
    \task $d$ và $d'$ song song với nhau.
    \task $d$ và $d'$ trùng nhau.
    \task $d$ và $d'$ cắt nhau.
    \task $d$ và $d'$ chéo nhau.
\end{tasks}
}{}

\vspace{0.15cm}

% Câu 32 (Câu 415 trong tài liệu - Vị trí tương đối 2 đường thẳng)
\questionLinesManual{0}{32}{%
Trong không gian với hệ toạ độ $Oxyz$, cho hai đường thẳng $d_1\colon \dfrac{x - 1}{1} = \dfrac{y + 3}{-2} = \dfrac{z + 3}{-3}$ và $d_2\colon \begin{cases} x = 3t \\ y = -1 + 2t \\ z = 0 \end{cases}$. Mệnh đề nào dưới đây đúng?
\begin{tasks}(2)
    \task $d_1$ song song với $d_2$.
    \task $d_1$ chéo $d_2$.
    \task $d_1$ cắt và vuông góc với $d_2$.
    \task $d_1$ cắt và không vuông góc với $d_2$.
\end{tasks}
}{}

\vspace{0.15cm}

% Câu 33 (Câu 417 trong tài liệu - Tham số k để 2 đường thẳng cắt nhau)
\questionLinesManual{0}{33}{%
Trong không gian với hệ tọa độ $Oxyz$, cho hai đường thẳng $d_1\colon \dfrac{x - 1}{1} = \dfrac{y - 2}{-2} = \dfrac{z - 3}{1}$ và $d_2\colon \begin{cases} x = 1 + kt \\ y = t \\ z = -1 + 2t \end{cases}$. Tìm giá trị của $k$ để $d_1$ cắt $d_2$.
\begin{tasks}(4)
    \task $k = 0$.
    \task $k = 1$.
    \task $k = -1$.
    \task $k = -\dfrac{1}{2}$.
\end{tasks}
}{}

\vspace{0.15cm}

% Câu 34 (Câu 418 trong tài liệu - Vị trí tương đối 2 đường thẳng)
\questionLinesManual{0}{34}{%
Trong không gian với hệ trục tọa độ $Oxyz$, cho hai đường thẳng $d_1\colon \dfrac{x - 1}{1} = \dfrac{y}{2} = \dfrac{z - 3}{3}$ và $d_2\colon \begin{cases} x = 2t \\ y = 1 + 4t \\ z = 2 + 6t \end{cases}$. Khẳng định nào sau đây là đúng?
\begin{tasks}(2)
    \task Hai đường thẳng $d_1, d_2$ song song với nhau.
    \task Hai đường thẳng $d_1, d_2$ trùng nhau.
    \task Hai đường thẳng $d_1, d_2$ cắt nhau.
    \task Hai đường thẳng $d_1, d_2$ chéo nhau.
\end{tasks}
}{}

\vspace{0.15cm}

% Câu 35 (Câu 420 trong tài liệu - Vị trí tương đối 2 đường thẳng)
\questionLinesManual{0}{35}{%
Trong không gian với hệ tọa độ $Oxyz$, cho hai đường thẳng $d_1\colon \begin{cases} x = t \\ y = 0 \\ z = 1 \end{cases}$ và $d_2\colon \begin{cases} x = 0 \\ y = t \\ z = 1 \end{cases}$. Khẳng định nào sau đây đúng?
\begin{tasks}(2)
    \task $d_1 // d_2$.
    \task $d_1$ và $d_2$ chéo nhau.
    \task $d_1$ và $d_2$ cắt nhau.
    \task $d_1 \equiv d_2$.
\end{tasks}
}{}

% =========================================================================
% ĐÁP ÁN BÀI TẬP VỀ NHÀ (35 CÂU MỚI 100%)
% =========================================================================
% 01. A | 08. A | 15. A | 22. B | 29. A
% 02. C | 09. A | 16. A | 23. C | 30. A
% 03. B | 10. B | 17. D | 24. D | 31. A
% 04. D | 11. B | 18. B | 25. A | 32. D
% 05. B | 12. D | 19. C | 26. B | 33. A
% 06. D | 13. D | 20. A | 27. C | 34. A
% 07. D | 14. D | 21. B | 28. A | 35. C

\end{document}